\unitopen{23}
\begin{candos}
\item mówić o planach (\de{Ich werde \dots}, \de{Ich habe vor \dots});
\item zbudować grzeczną prośbę z \de{würden}, \de{könnten}, \de{hätte gern}, \de{wäre};
\item powiedzieć, co chciałbyś robić w Szwajcarii;
\item odróżnić Präsens z określeniem czasu od \de{Futur I}.
\end{candos}

\section{Słownictwo}
\vocabtable{23}

\section{Przyszłość}
Najczęściej mówisz o przyszłości w Präsens z określeniem czasu: \de{Morgen fahre ich nach Zürich.} \de{Nächstes Jahr
ziehe ich um.} \emph{Futur I} (\de{werden} + bezokolicznik) podkreśla postanowienie lub przypuszczenie:
\de{Ich werde Deutsch lernen.} \de{Er wird wohl im Büro sein.}

\begin{center}\small
\begin{tabular}{@{}ll@{}}
\de{ich werde} & \de{wir werden}\\
\de{du wirst} & \de{ihr werdet}\\
\de{er/sie/es wird} & \de{sie/Sie werden}\\
\end{tabular}
\end{center}
Klamra jak zawsze: \de{Ich \textbf{werde} nächstes Jahr in Zürich \textbf{arbeiten}.}

\section{Uprzejmość: tryb przypuszczający w pięciu formach}
Grzeczna prośba często używa formy z \emph{trybu przypuszczającego} (\de{Konjunktiv II}, jednostka~29). Teraz uczysz się jej jako gotowych
zwrotów:

\begin{center}\small
\begin{tabular}{@{}lll@{}}
\toprule
\textsf{\textbf{Forma}} & \textsf{\textbf{Przykład}} & \textsf{\textbf{Funkcja}}\\
\midrule
\de{ich würde gern} & \de{Ich würde gern in Zürich arbeiten.} & życzenie\\
\de{ich hätte gern} & \de{Ich hätte gern einen Termin.} & zamawianie, prośba\\
\de{könnten Sie} & \de{Könnten Sie mir bitte helfen?} & uprzejma prośba\\
\de{wäre es möglich} & \de{Wäre es möglich, den Termin zu verschieben?} & bardzo uprzejma prośba\\
\de{ich wäre froh, wenn} & \de{Ich wäre froh, wenn Sie das prüfen könnten.} & miękka prośba\\
\bottomrule
\end{tabular}
\end{center}

\begin{falle}
Polski „chciałbym” to nie \de{Ich wollte}. \de{Ich wollte} to Präteritum („chciałem”). Grzeczne życzenie to
\de{Ich möchte} lub \de{Ich würde gern}. Drugi błąd: polskie „będę” z bezokolicznikiem tłumaczone zawsze przez
\de{werden}: po niemiecku Präsens z \de{morgen}, \de{nächste Woche} jest bardziej naturalne. Trzeci: \de{werden}
też znaczy „stawać się” (\de{Es wird kalt}), więc kontekst decyduje.
\end{falle}

\begin{szwajcaria}
Uprzejmość w szwajcarskim biurze jest ważniejsza niż bezpośredniość: \de{Könntest du kurz Zeit haben?} lub
\de{Hättest du einen Moment?} zamiast rozkazu. Szwajcarzy pytają często o plany: \de{Was sind Ihre Pläne?} --- na rozmowie
o pracę odpowiedź z \de{Ich würde gern bei Ihnen \dots{}} ma dobry efekt.
\end{szwajcaria}

\begin{praca}
\de{Ich würde gern an einem größeren Projekt arbeiten.} \de{Könnten Sie mir bitte das Dokument schicken?}
\de{Ich hätte gern Feedback zu meinem Pull Request.} \de{Wäre es möglich, das Meeting zu verschieben?}
\end{praca}

\section{Dialogi}
\textbf{Dialog 1 --- plany zawodowe}
\begin{dialog}
\turn{A}{Was haben Sie nächstes Jahr vor?}{Co ma Pan w planach na przyszły rok?}
\turn{B}{Ich werde Deutsch lernen und ich würde gern in der Schweiz arbeiten.}{Będę się uczył niemieckiego i chciałbym pracować w Szwajcarii.}
\turn{A}{Das wäre eine große Chance. Könnten Sie mir Ihren Lebenslauf schicken?}{To byłaby wielka szansa. Czy mógłby mi Pan wysłać CV?}
\turn{B}{Natürlich. Ich schicke ihn Ihnen morgen.}{Oczywiście. Wyślę je Panu jutro.}
\end{dialog}

\section{Ćwiczenia}
\begin{cwiczenia}
\exitem{Odmień \de{werden}: \de{ich \dots, du \dots, er \dots}}{\de{werde, wirst, wird}}
\exitem{Futur I: \de{Ich lerne Deutsch.} (pewność postanowienia)}{\de{Ich werde Deutsch lernen.}}
\exitem{Präsens z czasem: „Jutro jadę do Zurychu.”}{\de{Morgen fahre ich nach Zürich.}}
\exitem{Poproś grzecznie o pomoc (\de{könnten}).}{\de{Könnten Sie mir bitte helfen?}}
\exitem{Poproś o termin (\de{hätte gern}).}{\de{Ich hätte gern einen Termin.}}
\exitem{Wyraź życzenie pracy w Szwajcarii (\de{würde gern}).}{\de{Ich würde gern in der Schweiz arbeiten.}}
\exitem{Bardzo uprzejma prośba o przesunięcie (\de{wäre es möglich}).}{\de{Wäre es möglich, den Termin zu verschieben?}}
\exitem{Popraw: \de{Ich wollte gern einen Kaffee.} (chodzi o życzenie teraz)}{\de{Ich möchte gern einen Kaffee.} lub \de{Ich hätte gern \dots}}
\exitem{Wskaż Futur: \de{Er wird kommen.} czy \de{Es wird kalt.}?}{\de{Er wird kommen.} (Futur); \de{Es wird kalt} = „robi się zimno”}
\exitem{Przetłumacz: „Mam nadzieję, że znajdę pracę.”}{\de{Ich hoffe, dass ich eine Stelle finde.}}
\exitem{Przetłumacz: „Chciałbym się nauczyć niemieckiego.”}{\de{Ich würde gern Deutsch lernen.}}
\exitem{Przetłumacz: „Czy mógłby Pan powtórzyć?”}{\de{Könnten Sie das bitte wiederholen?}}
\end{cwiczenia}

\section{Trening słownictwa (generowany)}
\drill{23}

\weekend{Opisz w dziesięciu zdaniach swoje plany na najbliższy rok i pięciu lat: kariera, język, miejsce zamieszkania.
Użyj \de{werden}, \de{würde gern}, \de{hoffen, dass}. Potem nagraj trzy uprzejme prośby do przełożonego.}
